\section{Conclusion}
\label{sec:conclusion}

We asked whether the Avellaneda--Stoikov model survives contact with the National Stock Exchange of
India, and answered in three steps on a liquid and an illiquid stock.

\begin{enumerate}
\item \textbf{Assumptions.} The foundation holds: fill intensity decays exponentially in the
  quote distance, with $R^2$ above 0.98 for both stocks and both sides of the book. The description
  of the mid-price bends, mostly through the bid--ask bounce of the illiquid stock, and constant
  intraday volatility breaks.
\item \textbf{Risk.} The model's inventory control survives almost untouched. On held-out real
  minutes it cuts the standard deviation of daily PnL by a factor of \RELIANCEReplaySdFactor{}
  (RELIANCE) and \KSCLReplaySdFactor{} (KSCL), about the same as it does in its own simulated world.
\item \textbf{Profit.} The model's profit does not survive. Adverse selection, which the model assumes
  away, absorbs the whole spread income, and the model's spread is about three times wider than
  the market's, as one expects of a single dealer in a competitive order book.
\end{enumerate}

\noindent The practical reading is simple. On NSE, the Avellaneda--Stoikov rule is a sound way for a
liquidity provider to \emph{manage inventory}. It is not, on its own, a way to \emph{make money}. To
do that a market maker has to price the information in order flow, which is the problem the
information-based models of microstructure were built to solve.
